An image search system should reject wrong descriptions while recognizing different correct descriptions. Labels can mark a caption as supported, contradicted, or unresolved by an image. We ask what each use of these labels adds to retrieval. Twelve training rules use one pretrained model's fixed features, with matched caption pools and tuning budgets. Compared with training on the original captions, adding supported captions improves the annotated statement test by \PositiveRelationGain{} percentage points. It gains only \PositiveSugarGain{} points on edited captions and loses \PositiveRetrievalLoss{} points when retrieving captions for images. Gains concentrate on added objects and attributes; accuracy falls on replacement and swap edits. Extra penalties for contradicted captions and omitting unresolved pairs add no clear further gain. Uniformly reducing the influence of competing captions preserves more retrieval accuracy, but scores lower on the annotated task. Relation annotation must earn its value through the matches a search system returns.
